% Written by make_book from the data (band 1500-1599-average, two books, min-games 500, top 15, endgames 15, always pawns v pawns; B v B opposite, program be90d599; check 0ad09bf15ff6)
% If you edit this file, make_book keeps your version.
\chapter{How Often Each Endgame Comes Up}\label{ch:meet}
Before choosing what to study, it helps to know what you will actually meet over the board. This chapter counts, in all 22,946 games of players rated 1500–1599 (average of both players), how often each endgame came up. These are numbers from the games, not from the positions in this book. Only the endgames this book covers are shown, the same ones in all three charts of this part.

Each bar is one material balance, whoever has which side: “Rook vs rook” counts every game in which each side had a rook (and pawns). A game that passes through several endgames counts once for each of them.

\begin{center}
\begin{tikzpicture}
\begin{axis}[xbar,width=6.5cm,height=7.9cm,xmin=0,ytick={0,...,7},yticklabels={{Bishop and knight vs bishop and knight},{Queen and bishop vs queen and knight},{Rook and bishop vs two rooks},{Rook and knight vs rook and knight},{Queen and rook vs queen and rook},{Two rooks vs two rooks},{Rook and bishop vs rook and bishop},{Rook and bishop vs rook and knight}},ymin=-0.5,ymax=7.5,enlarge y limits=false,bar width=7pt,xmajorgrids,grid style={gray!20},axis x line*=bottom,axis y line*=left,tick label style={font=\footnotesize},clip=false]
\addplot[fill=bar,draw=none] coordinates {(1.329,0) (1.469,1) (1.691,2) (3.229,3) (4.284,4) (4.707,5) (4.707,6) (6.319,7)};
\node[anchor=west,font=\footnotesize,text=muted] at (axis description cs:1.02,0.0625) {1.3 of 100 games};
\node[anchor=west,font=\footnotesize,text=muted] at (axis description cs:1.02,0.1875) {1.5 of 100 games};
\node[anchor=west,font=\footnotesize,text=muted] at (axis description cs:1.02,0.3125) {1.7 of 100 games};
\node[anchor=west,font=\footnotesize,text=muted] at (axis description cs:1.02,0.4375) {3.2 of 100 games};
\node[anchor=west,font=\footnotesize,text=muted] at (axis description cs:1.02,0.5625) {4.3 of 100 games};
\node[anchor=west,font=\footnotesize,text=muted] at (axis description cs:1.02,0.6875) {4.7 of 100 games};
\node[anchor=west,font=\footnotesize,text=muted] at (axis description cs:1.02,0.8125) {4.7 of 100 games};
\node[anchor=west,font=\footnotesize,text=muted] at (axis description cs:1.02,0.9375) {6.3 of 100 games};
\end{axis}
\end{tikzpicture}
\end{center}
\begin{keypoints}
\item Rook and bishop vs rook and knight is the endgame you will meet most: it came up in 6 of every 100 games.
\item Next come rook and bishop vs rook and bishop (5), two rooks vs two rooks (5).
\end{keypoints}
\chapter{How Often Each Endgame Goes Wrong}\label{ch:go-wrong}
Some endgames are harder than others. This chapter asks: once an endgame is on the board, how often does a player throw away half a point or more? Again, these are numbers from all the games.

A blunder here is a move that turns a win into a draw or a loss, or a draw into a loss, while a better move was available. The bar shows how many such moves were made per 100 games that reached the endgame.

\begin{center}
\begin{tikzpicture}
\begin{axis}[xbar,width=6.5cm,height=7.9cm,xmin=0,ytick={0,...,7},yticklabels={{Rook and bishop vs two rooks},{Queen and rook vs queen and rook},{Two rooks vs two rooks},{Bishop and knight vs bishop and knight},{Rook and bishop vs rook and knight},{Rook and bishop vs rook and bishop},{Rook and knight vs rook and knight},{Queen and bishop vs queen and knight}},ymin=-0.5,ymax=7.5,enlarge y limits=false,bar width=7pt,xmajorgrids,grid style={gray!20},axis x line*=bottom,axis y line*=left,tick label style={font=\footnotesize},clip=false]
\addplot[fill=bar,draw=none] coordinates {(66.753,0) (87.284,1) (92.037,2) (99.016,3) (99.448,4) (105.926,5) (108.502,6) (110.386,7)};
\node[anchor=west,font=\footnotesize,text=muted] at (axis description cs:1.02,0.0625) {67 per 100 games};
\node[anchor=west,font=\footnotesize,text=muted] at (axis description cs:1.02,0.1875) {87 per 100 games};
\node[anchor=west,font=\footnotesize,text=muted] at (axis description cs:1.02,0.3125) {92 per 100 games};
\node[anchor=west,font=\footnotesize,text=muted] at (axis description cs:1.02,0.4375) {99 per 100 games};
\node[anchor=west,font=\footnotesize,text=muted] at (axis description cs:1.02,0.5625) {99 per 100 games};
\node[anchor=west,font=\footnotesize,text=muted] at (axis description cs:1.02,0.6875) {106 per 100 games};
\node[anchor=west,font=\footnotesize,text=muted] at (axis description cs:1.02,0.8125) {109 per 100 games};
\node[anchor=west,font=\footnotesize,text=muted] at (axis description cs:1.02,0.9375) {110 per 100 games};
\end{axis}
\end{tikzpicture}
\end{center}
\begin{keypoints}
\item Queen and bishop vs queen and knight is the hardest to handle: 110 blunders per 100 games that reach it.
\item Endgames that are hard to play are not always the ones you meet often. \xref{Chapter}{ch:where-to-start} puts both together.
\end{keypoints}
\chapter{Where to Start}\label{ch:where-to-start}
An endgame deserves your study time if it comes up often and goes wrong often. This chapter multiplies the two: how many blunders each endgame causes per 100 games played, whether the game reached that endgame or not.

Read the chart from the top. The endgames there cost players the most half points, so they come first in the rest of the book.

\begin{center}
\begin{tikzpicture}
\begin{axis}[xbar,width=6.5cm,height=7.9cm,xmin=0,ytick={0,...,7},yticklabels={{Rook and bishop vs two rooks},{Bishop and knight vs bishop and knight},{Queen and bishop vs queen and knight},{Rook and knight vs rook and knight},{Queen and rook vs queen and rook},{Two rooks vs two rooks},{Rook and bishop vs rook and bishop},{Rook and bishop vs rook and knight}},ymin=-0.5,ymax=7.5,enlarge y limits=false,bar width=7pt,xmajorgrids,grid style={gray!20},axis x line*=bottom,axis y line*=left,tick label style={font=\footnotesize},clip=false]
\addplot[fill=bar,draw=none] coordinates {(1.129,0) (1.316,1) (1.621,2) (3.504,3) (3.739,4) (4.332,5) (4.986,6) (6.284,7)};
\node[anchor=west,font=\footnotesize,text=muted] at (axis description cs:1.02,0.0625) {1.1 per 100 games};
\node[anchor=west,font=\footnotesize,text=muted] at (axis description cs:1.02,0.1875) {1.3 per 100 games};
\node[anchor=west,font=\footnotesize,text=muted] at (axis description cs:1.02,0.3125) {1.6 per 100 games};
\node[anchor=west,font=\footnotesize,text=muted] at (axis description cs:1.02,0.4375) {3.5 per 100 games};
\node[anchor=west,font=\footnotesize,text=muted] at (axis description cs:1.02,0.5625) {3.7 per 100 games};
\node[anchor=west,font=\footnotesize,text=muted] at (axis description cs:1.02,0.6875) {4.3 per 100 games};
\node[anchor=west,font=\footnotesize,text=muted] at (axis description cs:1.02,0.8125) {5.0 per 100 games};
\node[anchor=west,font=\footnotesize,text=muted] at (axis description cs:1.02,0.9375) {6.3 per 100 games};
\end{axis}
\end{tikzpicture}
\end{center}
\begin{keypoints}
\item Study in this order: rook and bishop vs rook and knight, then rook and bishop vs rook and bishop, then two rooks vs two rooks.
\item Rook and bishop vs rook and knight alone causes 6 blunders per 100 games played.
\item This book does not cover every endgame, and that is on purpose. It covers the 8 endgames in the chart, the ones that cost the most half points. The others, such as rook vs rook and bishop, come up rarely or go wrong seldom: the time you would spend on them is better spent here.
\end{keypoints}
